%% Modelo PIBIC
%% Modificado por: Juan Sebastian Toquica Arenas

\input{preamble}

\begin{document}

% Remove page number
\clearpage
\thispagestyle{empty}

\begin{bfseries}
\begin{center}

\includegraphics[scale=0.3]{ufc-cor.png} \\
\vspace{0.5cm} 
UNIVERSIDADE FEDERAL DO CEARÁ \\
\vspace{0.2cm} 
%CENTRO DE CIÊNCIAS \\
%\vspace{4pt} 
%DEPARTAMENTO DE FÍSICA \\
%\vspace{4pt}
CAMPUS ITAPAJÉ \\
\vspace{0.2cm}
JARDINS DE ANITTA \\
\vspace{4cm}
CHAMADA EDITAL PIBIC 2023/2024 \\


\vspace*{0.8cm}
\textbf{\MakeUppercase{Titulo projeto PIBIC}} \\
\vspace*{3cm}
\textbf{\MakeUppercase{prof. Juan Sebastian Toquica Arenas}}\\
\vspace*{7cm}
\textbf{\MakeUppercase{Itapajé, Ceará, abril de 2023}}

\end{center}

\end{bfseries}
\newpage

\section{\MakeUppercase{resumo}} 

\textcolor{red}{mínimo de 10 caracteres e máximo de 2000 caracteres}


\textbf{palavras chave:} chave1; chave2; chave3; chave4; chave5


\section{\MakeUppercase{introdução}} 

\textcolor{red}{(mínimo de 10 caracteres e máximo de 12300 caracteres)}

Início da parte textual do trabalho. Tem como finalidade dar ao leitor uma visão concisa do tema da prática, redigindo um texto de sua própria autoria sobre a fundamentação teórica envolvida na prática em questão. Aqui se pretende que o aluno exercite a habilidade de escrever em linguagem científica. Para redigir a introdução recomenda-se que o aluno pesquise em manuais, livros-texto e/ou Internet.

Todos os autores citados devem ter a referencia incluída em lista no final no trabalho. Não será tolerada a cópia total ou parcial sem a devida referência. A reprodução de figuras obtidas de outras fontes é permitida desde que seja citada a fonte, devidamente, no texto e nas referências.

Todo texto deve ser digitado em fonte Times New Roman, tamanho 12, inclusive a capa, com exceção das citações com mais de três linhas, notas de rodapé, paginação, legendas e fontes das ilustrações e das tabelas, que devem ser em fonte Times New Roman tamanho 10 e espaçamento entre linhas de 1,5. Com Margens: Superior e esquerda em três centímetros; inferior e direita em dois centímetros. O texto deve ser justificado, exceto as referências, no final do trabalho, que devem ser alinhadas a esquerda.


\jump

Para Siss (2012) as políticas de ação afirmativas constituem políticas públicas, estatais e de caráter compulsório, elaboradas e implementadas pelo Estado, ou seja, é o Estado em ação.

Segundo Bastos e Keller (2006, p. 38), “A leitura é um processo que envolve algumas habilidades, entre as quais a interpretação do texto e a sua compreensão.”

As organizações testemunharam uma redução da validade de seu conhecimento durante este período e começaram a perceber que já não era possível confiar em Instituições de Ensino Superior para desenvolver a sua mão de obra (TARAPANOFF, 2006).

O discurso jurídico, que hoje se apresenta com um novo perfil, dispõe de um acervo variado de opções para ser construído, pois, “[...] agrega valores, impõe condutas, conduz instituições, movimenta riquezas, opta por visões de mundo e, portanto, sustenta uma ideologia.” (BITTAR, 2001, p. 181).

\begin{leftquote}[]
A citação direta com mais de três linhas deve ser destacada com recuo de 4 cm da margem esquerda, em letra menor do que a utilizada no texto (tamanho 10), sem as aspas e com espaçamento simples entrelinhas. A citação deverá ser separada do texto que a precede e a sucede [...] (UNIVERSIDADE FEDERAL DO CEARÁ, 2013, p. 96).
\end{leftquote}

\begin{leftquote}[De acordo com Mueller e Perucchi (2014, p. 16)]
A expressão ciência e tecnologia e inovação, comumente abreviada para CT\&I, reflete a interação e a interdependência entre essas áreas. Pela sua natureza, o conhecimento científico, tecnológico e a inovação dependem da comunicação, e a comunicação desses conhecimentos é objeto de estudo da ciência da informação.
\end{leftquote}

\jump

As ilustrações (fotografias, gráficos, mapas, plantas, quadros) e tabelas devem ser citados e inseridos o mais próximo possível do trecho a que se referem.

\jump

% http://pgfplots.net/tikz/examples/bar-plot/ (the "right" way to do it)
\begin{container}{10cm}{Gráfico 1 – Distribuição dos documentos analisados por programa de pós-graduação}{Fonte: Elaborado pelo autor.}

\includegraphics[width=10cm]{graph.png}

\end{container}

\begin{container}{10cm}{Figura 1 – Organização do conhecimento/Representação do
conhecimento, Organização da informação/Representação da informação}{Fonte: Lara e Smit (2010, com adaptações).}

\includegraphics[width=10cm]{image.png}

\end{container}

\jump

\begin{container}{10cm}{Tabela 1 – Distribuição dos documentos analisados por programa de pós-graduação}{Fonte: elaborada pelo autor.}
\begin{center}
\begin{tabular}{ | c | c | c | c | }
    \hline
    Programas de pós-graduação & Teses & Dissertações & Total\\
    \hline
    \hline
    Cirurgia & 1 & 1 & 2 \\
    \hline
    Enfermagem & 4 & 4 & 8 \\
    \hline
    Engenharia Civil & 2 & 8 & 10 \\
    \hline
    Farmacologia & 8 & 6 & 14 \\
    \hline
    Física & 3 & 6 & 9 \\
    \hline
    Química Inorgânica & 4 & 1 & 5 \\
    \hline
    Total & 22 & 26 & 48 \\
    \hline
\end{tabular} 
\end{center}
\end{container}

\section{\MakeUppercase{Perguntas de partida}}

\textcolor{red}{Preenchimento opcional - máximo de 2000 caracteres}

Nesta parte o aluno deve descrever detalhadamente os procedimentos realizados (em primeira pessoa); apresentar os resultados das medidas (tabelas, gráficos, cálculos, etc). Atenção: anotar sempre as unidades.


\section{\MakeUppercase{Hipóteses}}

\textcolor{red}{Preenchimento opcional - máximo de 2000 caracteres}
No questionário devem constar as perguntas e as respostas de acordo com o roteiro de prática do ano em curso. 


\section{\MakeUppercase{Objetivos}}

\textcolor{red}{mínimo de 10 caracteres e máximo de 4100 caracteres}

Parte final do texto na qual se apresentam as conclusões apoiadas no desenvolvimento do assunto. É a recapitulação sintética dos resultados obtidos. 

Na conclusão o aluno deve apresentar uma discussão sobre os resultados obtidos em função dos objetivos do experimento. Cabe também comentar sobre as possíveis fontes de erros. \cite{alvares2017retrofitting}

\cite{hosseinzadeh2021improving}

\section{\MakeUppercase{Metodologia}}

\textcolor{red}{mínimo de 10 caracteres e máximo de 16400 caracteres}

\section{\MakeUppercase{Dados preliminares}}

\textcolor{red}{Preenchimento opcional - máximo de 8200 caracteres}

\newpage

\section{\MakeUppercase{Referências}}

\renewcommand\refname{}
\bibliographystyle{abntex2-num}
\bibliography{bibfile}


\newpage

\section{\MakeUppercase{Plano de atividades}}


\end{document}
